% Felipe Muñoz Casasola
% This is free and unencumbered software released into the public domain.

% Anyone is free to copy, modify, publish, use, compile, sell, or
% distribute this software, either in source code form or as a compiled
% binary, for any purpose, commercial or non-commercial, and by any
% means.

% In jurisdictions that recognize copyright laws, the author or authors
% of this software dedicate any and all copyright interest in the
% software to the public domain. We make this dedication for the benefit
% of the public at large and to the detriment of our heirs and
% successors. We intend this dedication to be an overt act of
% relinquishment in perpetuity of all present and future rights to this
% software under copyright law.

% THE SOFTWARE IS PROVIDED "AS IS", WITHOUT WARRANTY OF ANY KIND,
% EXPRESS OR IMPLIED, INCLUDING BUT NOT LIMITED TO THE WARRANTIES OF
% MERCHANTABILITY, FITNESS FOR A PARTICULAR PURPOSE AND NONINFRINGEMENT.
% IN NO EVENT SHALL THE AUTHORS BE LIABLE FOR ANY CLAIM, DAMAGES OR
% OTHER LIABILITY, WHETHER IN AN ACTION OF CONTRACT, TORT OR OTHERWISE,
% ARISING FROM, OUT OF OR IN CONNECTION WITH THE SOFTWARE OR THE USE OR
% OTHER DEALINGS IN THE SOFTWARE.

% For more information, please refer to <https://unlicense.org/>

\documentclass{article}

\usepackage[spanish]{babel}
\usepackage[pdfusetitle]{hyperref}
\usepackage{microtype}

\title{Gestión de la página web}
\author{Felipe Muñoz Casasola}
\date{13 de febrero de 2026}

\begin{document}
    Estoy siguiendo las instrucciones de
    \href{https://mullvad.net/en/help/how-blog-anonymously}{https://mullvad.net/en/help/how-blog-anonymously},
    archivado
    \href{https://web.archive.org/web/20260102025718/https://mullvad.net/en/help/how-blog-anonymously}{aquí}.
    Voy a usar una máquina virtual llamada Página-web que tiene por sistema
    operativo Ubuntu. La contraseña del único usuario que hay es el mismo
    nombre de usuario (felmuncas) . Hay que falsificar los datos que se le den
    al sistema operativo.

    Siempre que vaya a escribir algún documento que pueda ser público lo haré a
    través de este sistema operativo, donde utilizaré una instalación local de
    \LaTeX{}. De esta forma evito que se pueda reconocer el estilo de
    indentación u otro estilo que tenga personalizado en mi configuración que
    sea reconocible. Tampoco tengo en público la configuración de los «snippets»
    de VimTex, pues esto al final refleja la estructura de todos los documentos
    de \LaTeX{} que hago, y si luego publico por otro lado documentos en
    \LaTeX{} con la misma estructura se me puede identificar.

    También tendré activada la VPN de Proton, y programaré en \LaTeX{} en
    Codium (para evitar la tentación, si alguna vez se me olvida, de
    usar vimtex copiando la configuración que he preferido mantener en privado).

    La idea es que los documentos PDF finales sean compilados por la máquina virtual.
    Pero en un futuro (tengo que arreglarlo) deberé incrustar los archivos de configuración
    de mi ordenador con minted, de forma que se cargarán de mi ordenador en los
    documentos de \LaTeX{} en tiempo de ejecución. Para esto parece ser que
    actualmente \LaTeX{} usa una carpeta \texttt{\_minted}, que parece que la
    reutiliza el compilador cuando compilo a través de la máquina virtual.
    Por tanto, deberé primero compilar (usando el \texttt{Makefile} del
    proyecto) en mi ordenador principal, para que se creen las carpetas
    \texttt{\_minted} con los archivos de configuración incrustados. Y que
    entonces en la máquina virtual al compilar se recuperen esos textos.

    Por tanto, para actualizar la página web, tengo que modificarla en la
    máquina virtual. Posteriormente llevármela a mi ordenador principal,
    ejecutar make, volvérmelo a llevar a la máquina virtual y una vez allí
    hacer \texttt{make pre}, para observar cómo quedaría la página con los
    nuevos cambios. Y una vez comprobado que funciona bien, hacer
    \texttt{make pub} para publicar los cambios recompilando primero.21
    Espero que con esto se puedan aprovechar los archivos de \texttt{minted}
    y no me dé errores de compilación.

    Tampoco debo acceder nunca a mí página web desde fuera de la máquina
    virtual para que no se me pueda relacionar en absoluto. Si tuviera que
    consultar las guías que guardo, simplemente uso el USB en el que se están
    guardando. Pues como tiene adaptador de USB tipo C también lo puedo
    conectar al teléfono si me hiciera falta.
\end{document}
